
\documentclass[11pt,a4paper]{article}
\usepackage[utf8]{inputenc}
\usepackage[T1]{fontenc}
\usepackage[ngerman]{babel}
\usepackage[margin=2.3cm]{geometry}
\usepackage{booktabs}
\usepackage{amsmath}
\usepackage{lmodern}
\usepackage[protrusion=true,expansion=false]{microtype}
\usepackage[hidelinks]{hyperref}
\usepackage{parskip}
\usepackage{fancyhdr}
\pagestyle{fancy}
\fancyhf{}
\renewcommand{\headrulewidth}{0pt}
\renewcommand{\footrulewidth}{0.4pt}
\fancyfoot[L]{\footnotesize\href{https://github.com/gerontec/blastDiamond}%
{\texttt{github.com/gerontec/blastDiamond}}}
\fancyfoot[R]{\footnotesize\thepage}
\fancypagestyle{plain}{\fancyhf{}%
  \renewcommand{\headrulewidth}{0pt}\renewcommand{\footrulewidth}{0.4pt}%
  \fancyfoot[L]{\footnotesize\href{https://github.com/gerontec/blastDiamond}%
  {\texttt{github.com/gerontec/blastDiamond}}}%
  \fancyfoot[R]{\footnotesize\thepage}}
\usepackage{listings}
\lstset{basicstyle=\ttfamily\scriptsize, breaklines=true,
        breakatwhitespace=false, columns=fullflexible,
        frame=single, framesep=4pt, xleftmargin=0pt,
        postbreak=\mbox{\textcolor{gray}{$\hookrightarrow$}\space}}
\usepackage{xcolor}
\setlength{\emergencystretch}{3em}
\sloppy

\title{\textbf{+1-Frameshift bei m1$\Psi$-mRNA}\\[2mm]
\large Was das Nature-Paper belegt und was die nr-Datenbank dazu hergibt}
\author{Auswertung auf dell-3660, NCBI-nr in sechs DIAMOND-Teil-Datenbanken}
\date{@@datum@@}

\begin{document}
\maketitle

\section*{Die Fragen und die kurzen Antworten}

Anlass war ein Zeitungstext über „gestörte Translation``. Er bezieht sich auf
Mulroney et al., \emph{Nature} \textbf{625}, 189--194 (2024),
\emph{N1-methylpseudouridylation of mRNA causes +1 ribosomal frameshifting}
(DOI \href{https://doi.org/10.1038/s41586-023-06800-3}{10.1038/s41586-023-06800-3},
PMID 38057663).

\begin{enumerate}
\item \textbf{Gibt es in der BLAST-DB schon Proben solcher Proteine?}
      Nein. Weder vorher zu erwarten noch nachträglich gefunden.
\item \textbf{Nennt das Paper eine solche Sequenz?} Ja, zwei Sätze:
      @@n_fs@@ vorhergesagte Peptide in Extended Data Table~3 und 9
      massenspektrometrisch belegte in Extended Data Table~1.
\item \textbf{Ein humanes Protein wird nicht als human erfasst?} Richtig, und
      das ist kein Erfassungsfehler: die Taxon-Angabe in nr beschreibt die
      Herkunft der kodierenden Sequenz, nicht die Zelle, die translatiert hat.
\item \textbf{In welcher Teil-DB findet sich etwas?} Siehe
      Abschnitt~\ref{sec:ergebnis}; die Positivkontrolle trifft, die
      Impfstoff-Peptide nicht.
\item \textbf{Wurde der Abbau über Entzündungsmarker untersucht?} Nein, von
      niemandem.
\item \textbf{Wie wurden die Proben ausgewählt, und wie hoch war die
      Trefferquote?} Nicht randomisiert, Gelegenheitsstichprobe aus vier
      Biobanken; Trefferquote 5 von 21 Geimpften, also rund 24\,\% -- nicht
      1 von 1000.
\item \textbf{Was ist das ähnlichste Protein in der ganzen nr?}
      \texttt{AAG23172.1}, ein menschlicher Eintrag, mit 88\,\% Identität über
      50 Reste -- und dieser eine Treffer stellt einiges richtig, siehe
      Abschnitt~\ref{sec:aehnlich}.
\end{enumerate}

\section{Was im Paper tatsächlich steht}

\subsection{Mechanismus}

N1-Methylpseudouridin (m1$\Psi$) verlangsamt die Translation und lässt das
Ribosom an \emph{slippery sites} um eine Base nach vorn rutschen. Die Autoren
definieren diese Stellen über
\[
\text{m1}\Psi\;\text{m1}\Psi\;\text{m1}\Psi\;X,
\qquad X \in \{\text{m1}\Psi,\ \text{C}\},
\]
wobei $X$ das \emph{erste} Nukleotid des unmittelbar folgenden Codons ist --
faktisch \texttt{UUU} gefolgt von \texttt{U} oder \texttt{C}. Im Fluc-Reporter
liegen drei solche Stellen (A, B, C). A liegt vor einem +1-Stoppcodon und trägt
nichts bei; die Mutanten U*187C und U*208C (Stellen B und C) senken das
Frameshifting deutlich, die Doppelmutante auf nicht mehr nachweisbar.
\textbf{Sechs} identische Stellen verteilen sich über die BNT162b2-Spike-CDS
(Extended Data Fig.~6).

Rate: rund 8\,\% im Reporter; die Pressemitteilung aus Cambridge nennt etwa
10\,\%. Meist bricht die Translation nach dem Rutsch ab. Selten läuft sie im
+1-Rahmen weiter und erzeugt ein chimäres Protein: vorn korrekt, hinter der
Rutschstelle sinnlos. Das Junction-Peptid über die Rahmengrenze wurde per
LC--MS/MS direkt nachgewiesen (Fig.~3b), dazu neun Peptide aus dem +1-Rahmen
(Extended Data Table~1).

\subsection{Nachweis am Menschen}

IFN$\gamma$-ELISpot: PBMC von 21 BNT162b2-Geimpften gegen 20
ChAdOx1-Geimpfte, stimuliert mit den +1FS-Peptiden. Signifikant erhöhte Antwort
nur nach BNT162b2 ($P = 0{,}0233$); bei Mäusen deutlicher ($P = 0{,}0005$ gegen
unbehandelt). ChAdOx1 nutzt keine m1$\Psi$-mRNA und dient als Negativkontrolle.
Auf In-Frame-Spike reagierten beide Gruppen.

\subsection{Wie die Proben ausgewählt wurden}

Nicht randomisiert. Das Paper enthält weder einen Kohortenabschnitt noch
Ein- oder Ausschlusskriterien; die Proben stammen als Gelegenheitsstichprobe
aus vier bereits bestehenden britischen Sammlungen:

\begin{itemize}
\item NIHR National BioResource, Research Tissue Bank, REC 17/EE/0025
\item PITCH, Substudie von SIREN (IRAS 284460, REC 20/SC/0230,
      ISRCTN11041050) -- SIREN ist eine Kohorte von Beschäftigten im
      Gesundheitswesen, also keine Bevölkerungsstichprobe
\item Liverpool, \emph{Human immune responses to acute virus infections}
      (16/NW/0170)
\item Oxford, GI Biobank Study (16/YH/0247)
\end{itemize}

Die einzige genannte Auswahleigenschaft der Teilnehmer ist, dass
\emph{``none of whom reported undue effects as a result of vaccination''}.
Alter und Geschlecht stehen nur nachträglich in Supplementary Table~1, nicht
als Kriterium. Für die Tierversuche sagt das Paper es ausdrücklich:
\emph{``Mice were not randomly assigned to groups and experimenters were not
blinded to experiments''} ($n = 8$ je Gruppe).

\subsection{Trefferquote: etwa jeder Vierte}

Die Einzelwerte stehen in Extended Data Fig.~2: die Sample-IDs kodieren das
Responder-Flag und die SFU-Zahl, Schwelle ist ein ELISpot-Count
$> 40$\,SFU/Million. Ausgezählt ergibt das fünf Responder:

\begin{center}
\begin{tabular}{lr}
\toprule
Responder & SFU/Mio. \\
\midrule
Donor-11 & 155 \\
Donor-21 & 105 \\
Donor-19 & 51{,}7 \\
Donor-25 & 43 \\
Donor-13 & 40 \\
\bottomrule
\end{tabular}
\end{center}

Dem stehen 15 Geimpfte unter der Schwelle gegenüber (Donor-32 bei 10, alle
übrigen bei 0). Das sind die 20 BNT162b2-Geimpften mit HLA-Typisierung; der
ELISpot-Vergleich nennt 21, einem fehlt die HLA-Bestimmung.

\[
\text{Trefferquote} \;=\; \frac{5}{21} \;\approx\; 24\,\%
\]

Zwei Einschränkungen sind der Abbildung direkt zu entnehmen. Donor-13 liegt mit
exakt 40 \emph{nicht} über der eigenen Schwelle „$> 40$``, wird aber als
Responder geführt. Und die Signifikanz der Humandaten ($P = 0{,}0233$) stammt
aus einem \textbf{einseitigen} Welch-t-Test; bei fünf Respondern unter 21, von
denen zwei knapp an der Schwelle liegen, ist das ein schwaches Ergebnis. Die
Mausdaten sind mit $P = 0{,}0005$ bzw.\ $P = 0{,}001$ deutlich robuster.

Davon zu trennen ist die \emph{molekulare} Rutschrate: rund 8\,\% im
Fluc-Reporter, in der Cambridge-Pressemitteilung mit etwa 10\,\% angegeben --
also ungefähr 1 von 12, nicht 1 von 1000. Die beiden Zahlen messen
Verschiedenes: 8\,\% ist Biochemie im Reagenzglas, 24\,\% ist der Anteil
Geimpfter, bei denen später eine T-Zell-Antwort gegen diese Produkte
nachweisbar war. In der Regel bricht die Translation nach dem Rutsch ohnehin
ab; nur ein Teil läuft im +1-Rahmen weiter und ergibt ein vollständiges
Fehlprodukt.

\subsection{Abbau und Entzündungsmarker: nicht untersucht}

Ein Volltext-Scan des Papers ergibt \textbf{null} Fundstellen für
\texttt{degrad*}, \texttt{proteasom*}, \texttt{inflammat*}, \texttt{half-life}
und \texttt{misfold*}. Es gibt keine Abbaukinetik, keine CRP- oder IL-6-Werte,
keinen Zeitverlauf.

Indirekt folgt aus der ELISpot-Antwort allerdings zwingend, dass die Produkte
\emph{überhaupt} abgebaut werden: MHC-I-Präsentation setzt proteasomalen Abbau
voraus. Dass T-Zellen reagieren, heißt also, die Produkte werden zerlegt und
präsentiert. Über die Geschwindigkeit sagt das nichts.

Die Autoren benennen die Lücke selbst: \emph{``the evolution of antibody and
cytolytic T cell responses against +1 frameshifted spike variants and peptides
has not been systematically evaluated in humans''}, und zur Risikofrage:
\emph{``there is no evidence that frameshifted products in humans generated
from BNT162b2 vaccination are associated with adverse outcomes''}.

Eine PubMed-Suche \texttt{frameshift AND (BNT162b2 OR N1-methylpseudouridine)}
liefert genau einen Treffer: das Paper selbst. Eine Folgestudie mit
Entzündungsmarkern existiert nicht. Was als Folgearbeit zirkuliert -- Boros et
al., \emph{Pharmacol Res Perspect} \textbf{12}, e1218 (2024), PMID 38867495 --
ist ein Übersichtsartikel ohne eigene Messdaten zum Frameshift; die dort
referierten neun +1-Peptide sind Mulroneys Extended Data Table~1. Die dort
angeführte erhöhte FDG-Aufnahme im Herzmuskel bis 180 Tage ist Interpretation
der Autoren und unabhängig nicht bestätigt.

\section{Warum nr diese Proteine nicht führt, und erst recht nicht als human}

Die \texttt{taxid} in nr beschreibt die Herkunft der \emph{kodierenden
Sequenz}, nicht die Zelle, die das Protein gebaut hat. nr ist keine
Proteomvorhersage, sondern eine Sammlung eingereichter Records
(GenBank-CDS-Translationen, RefSeq, SwissProt, PDB, PIR, PRF); der Taxon stammt
aus dem \texttt{/organism}-Qualifier des Einreichers. Daraus folgt dreifach:

\begin{enumerate}
\item Kodiert werden die Produkte von der BNT162b2-mRNA, einem synthetischen,
      codon-optimierten Konstrukt. Dass menschliche Ribosomen translatieren,
      macht daraus kein humanes Gen -- so wie auch das In-Frame-Spike in nr
      unter SARS-CoV-2 bzw.\ \emph{synthetic construct} läuft, nicht unter
      \emph{Homo sapiens}.
\item Es ist ein \emph{Translations}fehler. Die Autoren haben per RNA-seq
      geprüft, dass weder DNA-Mutation noch Transkriptfehler vorliegen
      (Extended Data Table~2). RefSeq annotiert Gene und deren korrekte
      Leserahmen; für Produkte fehlerhafter Translation gibt es in nr keine
      Kategorie.
\item Niemand hat sie eingereicht. Hinterlegt sind nur MassIVE
      \texttt{MSV000093074} (MS-Rohdaten), GEO \texttt{GSE223044} (RNA-seq) und
      figshare \texttt{10.6084/m9.figshare.24271744}. Die vorhergesagten
      +1-Polypeptide existierten nur als lokale FASTA in der
      Proteome-Discoverer-Suche, ohne Accession.
\end{enumerate}

Was man \emph{im} humanen Taxon findet, ist das Gegenteil: echte
Frameshift-Produkte aus \emph{Mutationen}, weil dort das Gen selbst verschoben
ist und eingereicht wurde. Der dokumentierte Fall dazu im Repo ist UBB+1
(\texttt{examples/frameshift\_ubb1.md}): Ubiquitin 1--75 plus 20 fremde Reste,
in nr als \texttt{3K9O\_B} hinterlegt.

Die Autoren haben den Abgleich sogar selbst gemacht: Extended Data Fig.~4,
\emph{``NCBI BLASTP alignment of +1 translated products''}, BNT162b2-+1-Rahmen
gegen den +1-Rahmen des Wildtyp-S-Gens (NC\_045512.2). Ergebnis im Text: die
+1-Peptide folgen aus der Impfstoffsequenz, \emph{nicht} aus dem Wildvirus.

\section{Die Suche gegen die lokale nr}

\subsection{Abfragesatz}

Drei Klassen, unterscheidbar am FASTA-Präfix:

@@class_summary@@

Die Positivkontrolle ist der entscheidende Teil des Aufbaus: Firefly-Luciferase
steht in nr (UniProt \texttt{P08659}, 550\,aa). Werden ihre In-Frame-Peptide
\emph{nicht} gefunden, ist die Suche falsch konfiguriert und ein negatives
Ergebnis bei den Impfstoff-Peptiden wertlos.

Die Abschrift der 123 Peptide aus dem Tabellenbild wurde über die Kettenlogik
des Pools geprüft: aufeinander folgende Peptide überlappen, solange sie zum
selben +1-Abschnitt gehören. Es bleiben 41 Abbruchstellen, also 42
+1-Abschnitte, die jeweils an einem Stoppcodon enden; alle übrigen 81 Übergänge
passen buchstabengenau zusammen. Eine Inkonsistenz
(\texttt{PPDSPLCTPGTGSGSAI} gegen \texttt{LCTPPGTGSGSAIAW}, ein \texttt{P}
Unterschied) steht so in der Publikation und wurde nicht korrigiert.

\subsection{Parameter}

Sechs DIAMOND-Teil-DBs der NCBI-nr, zusammen rund 1,18\,Mrd.\ Sequenzen.
\texttt{-{}-dbsize} auf die Summe aller Teile gesetzt, damit die E-Werte denen
einer Suche gegen die ganze nr entsprechen. Sensitivität: DIAMOND-Standard, wie
in \texttt{pipeline/frameshift\_scan.py}, damit das Ergebnis mit dem
dokumentierten UBB+1-Lauf vergleichbar bleibt; \texttt{-{}-evalue 1000},
\texttt{-{}-max-target-seqs 50}, 20 Threads, \texttt{-{}-block-size 2}.

Die lockere E-Wert-Schwelle ist nötig, weil Peptide von 8--27\,aa niedrige
Scores erreichen. Exaktheit wird erst nachgelagert erzwungen: gezählt wird nur,
was \texttt{pident} $= 100$ \emph{und} \texttt{length} $=$ \texttt{qlen}
erfüllt, also das ganze Peptid buchstabengenau abdeckt.

Als Empfindlichkeitskontrolle lief p1\_mammalia zusätzlich mit
\texttt{-{}-ultra-sensitive -{}-masking 0 -{}-comp-based-stats 0}: 457\,s statt
18\,s, mit demselben Ergebnis für die Spike-Peptide. Die Standardeinstellung
übersieht hier also nichts.

\subsection{Laufzeiten}

@@runtime_table@@

Gesamtlaufzeit @@total_runtime@@ für @@n_all@@ Peptide gegen 1,18\,Mrd.\ Sequenzen.
Die Laufzeit folgt der Dateigröße der Teil-DB, nicht der Trefferzahl: DIAMOND
liest den Block-Index linear.

\subsection{Ergebnis}
\label{sec:ergebnis}

DIAMOND lieferte roh @@n_raw@@ HSPs. Nach dem Exaktheitsfilter bleiben
@@n_exact@@ Zeilen, die gegen \texttt{blast\_acc}, \texttt{blast\_seq} und
\texttt{blast\_taxon} in der wagodb verbunden wurden:

\footnotesize
@@meta_table@@
\normalsize

\subsection{Einordnung}

Drei Dinge sind daran bemerkenswert.

\textbf{Erstens} greift die Positivkontrolle: die In-Frame-Luciferase-Peptide
werden gefunden. Die Suche funktioniert, das negative Ergebnis bei den
Spike-Peptiden ist also belastbar.

\textbf{Zweitens} werden auch die MS-belegten „+1FS``-Peptide des Reporters
gefunden -- das ist kein Widerspruch, sondern Bauart. Im Fluc+1FS-Konstrukt ist
der C-Terminus der Luciferase \emph{absichtlich} in den +1-Rahmen gelegt; wenn
das Ribosom rutscht, entsteht echte Luciferase-Sequenz, und genau deshalb
kehrt die Enzymaktivität zurück. Diese Peptide sind Reporter-Artefakte, keine
Impfstoffprodukte.

\textbf{Drittens} trifft kein einziges der @@n_fs@@ Spike-Peptide. Bei
1,18\,Mrd.\ Sequenzen tritt nicht einmal ein \emph{zufälliger}
Volllängentreffer auf -- bei Peptiden ab etwa 12\,aa ist der Sequenzraum
($20^{12} \approx 4 \cdot 10^{15}$) dafür zu groß. Das bestätigt den Befund der
Autoren mit unabhängigen Daten und einem anderen Werkzeug: die +1-Produkte sind
Artefakte der codon-optimierten Impfstoffsequenz und haben in der Natur kein
Gegenstück.

Methodische Grenze: DIAMOND ist ein Seed-and-Extend-Verfahren und für sehr
kurze Queries nicht die empfindlichste Wahl. Für die kürzesten Peptide des
Pools (\texttt{WSQSCCVA} und \texttt{PQRSCLCP}, je 8\,aa;
\texttt{PRMCCTRTR}, 9\,aa) kann ein exakter Treffer theoretisch durch das
Seed-Raster fallen. Für die Fragestellung ist das unerheblich: ein
8-mer-Treffer wäre Zufall und kein Beleg für Existenz.

\section{Clusterlauf: gruppieren die Peptide mit etwas Menschlichem?}

Die exakte Suche beantwortet nur, ob eine Sequenz \emph{identisch} vorkommt.
Die schwächere Frage -- ob ein Peptid mit einem menschlichen Protein
\emph{verwandt} genug ist, um mit ihm in einer Gruppe zu landen -- beantwortet
ein Clusterlauf.

Das bricht bewusst die stehende Regel von
\texttt{pipeline/nr\_cluster\_daily.py} (``new sequences are compared only with
each other, not with the rest of nr''). Die Peptide stehen in keiner
nr-Ausgabe, der tägliche Lauf um 05:30 würde sie also nie anfassen. Der
Gegenpart wurde auf \emph{Homo sapiens} beschränkt (taxid 9606,
2\,285\,891 Sequenzen aus 3\,630\,866 Accessions), was den Lauf bezahlbar hält
und genau die interessierende Frage trifft.

Zwei Ebenen wie im Tageslauf, 90\,\% und darauf 50\,\% Identität. Ein
Parameter musste abweichen: der Tageslauf nutzt \texttt{-{}-cov-mode 1}
(Abdeckung des Targets), womit ein 15\,aa-Peptid niemals zu einem 550\,aa-Protein
gehören kann, weil 80\,\% des Proteins abgedeckt sein müssten. Hier steht
\texttt{-{}-cov-mode 2} (Abdeckung der kürzeren Sequenz), damit ein Peptid
überhaupt Mitglied werden \emph{kann}. Mit \texttt{-{}-cov-mode 1} wäre jedes
Peptid konstruktionsbedingt ein Einzelgänger -- ein Ergebnis ohne Aussage.

In \texttt{blast\_cluster} und \texttt{blast\_cluster\_rep} wurde nichts
geschrieben: diese Tabellen sind über \texttt{(oid, part)} verschlüsselt und
tragen die Semantik eines Tagesupdates. Dies hier ist eine Auswertung, ihr
Ergebnis liegt als TSV.

@@cluster_section@@

\section{Das ähnlichste Protein in der ganzen nr}
\label{sec:aehnlich}

Die Exaktsuche beantwortet nur ja/nein. Die naheliegende Anschlussfrage --
\emph{was ist das Ähnlichste, was nr überhaupt zu bieten hat} -- lässt sich aus
denselben Rohdaten beantworten, ohne neuen Lauf: bester HSP je Peptid, sortiert
nach Bitscore.

Von den 123 Spike-Peptiden haben \textbf{acht} überhaupt irgendeinen Treffer,
115 gar keinen. Und die fünf besten gehören zusammen:

\begin{center}
\footnotesize
\begin{tabular}{llrrrl}
\toprule
Peptid & Abdeckung & Ident. & Bits & E-Wert & Subject \\
\midrule
\texttt{fs116} & 17/18 & 94{,}1\,\% & 42{,}0 & 4{,}5e--03 & \texttt{AAG23172.1} \emph{Homo sapiens} \\
\texttt{fs119} & 17/18 & 100\,\% & 41{,}6 & 6{,}4e--03 & \texttt{AAG23172.1} \emph{Homo sapiens} \\
\texttt{fs118} & 17/17 & 88{,}2\,\% & 39{,}3 & 4{,}8e--02 & \texttt{AAG23172.1} \emph{Homo sapiens} \\
\texttt{fs117} & 16/17 & 93{,}8\,\% & 37{,}4 & 2{,}7e--01 & \texttt{AAG23172.1} \emph{Homo sapiens} \\
\texttt{fs120} & 15/18 & 100\,\% & 35{,}4 & 1{,}5e+00 & \texttt{AAG23172.1} \emph{Homo sapiens} \\
\midrule
\texttt{fs082} & 15/15 & 66{,}7\,\% & 28{,}9 & 5{,}0e+02 & \texttt{KAF9461180.1} \emph{Collybia nuda} \\
\texttt{fs025} & 18/18 & 50{,}0\,\% & 28{,}5 & 8{,}1e+02 & \texttt{CEM26774.1} \emph{Chromera velia} \\
\texttt{fs006} & 15/15 & 66{,}7\,\% & 28{,}1 & 9{,}2e+02 & \texttt{KAF2594081.1} \emph{Brassica cretica} \\
\bottomrule
\end{tabular}
\end{center}

Die unteren drei sind Rauschen: E-Werte von 500 bis 900, also reine
Zufallsähnlichkeit. Die oberen fünf sind es nicht. \texttt{fs116} bis
\texttt{fs120} sind fünf aufeinander folgende, überlappende Peptide des Pools;
aneinandergelegt ergeben sie einen zusammenhängenden Abschnitt von 50 Resten:

\begin{center}
\texttt{SSWYCMHAMLAAPFPSWVPRVSPDLGSQVCSHLHLPHSPPLLVPDTSQAR}
\end{center}

Gegen \texttt{AAG23172.1} (\emph{Homo sapiens}, 103\,aa, Gen \texttt{PP2150},
Titel schlicht ``unknown'') ausgerichtet, ohne Lücken:

\begin{lstlisting}
 +1                                       SSWYCMHAMLAAPFPSWVPRVSPD
                                           | |||||||||||||||| |||||
AAG   MGDQVPREQDLAKRSSPGPSPVLNMCNRQSAWATALSPWYCMHAMLAAPFPSWAPRVSPD

 +1   LGSQVCSHLHLPHSPPLLVPDTSQAR
        ||||||||||||||||||||||
AAG   PGSQVCSHLHLPHSPPLLVPDTSTPTWSSPIAHKRGGTRDELS
\end{lstlisting}

\textbf{44 von 50 Resten identisch, 88\,\%.} Bei 1,18\,Mrd.\ Sequenzen ist das
kein Zufall, und es widerspricht auf den ersten Blick allem Übrigen in diesem
Bericht.

\subsection{Auflösung: das ist gar kein Spike}

\texttt{AAG23172.1} stammt aus \texttt{AF269289.1}, einer 1035\,bp
``Homo sapiens PP2150 mRNA'', eingereicht am 17.05.2000 vom Shanghai Cancer
Institute, unveröffentlicht, Produkt ``unknown'', Vermerk \emph{``Method:
conceptual translation supplied by author''}. Das Protein ist also selbst nur
ein vom Einreicher geratener Leserahmen.

Entscheidend sind die Nukleotide. Die 150\,bp von \texttt{AF269289.1}, die
den übereinstimmenden Proteinabschnitt kodieren (nt 665--814), wurden gegen die
mRNA von \emph{TLE5} (früher \emph{AES}) geprüft:

\begin{lstlisting}
NM_198969.1  TLE5 transcript variant 1, 1754 nt
  laengste exakt uebereinstimmende Teilsequenz: 150 von 150 nt, ab Position 1216
\end{lstlisting}

Vollständig identisch, im 3'-untranslatierten Bereich von \emph{TLE5}. Und
genau dort liegt die Erklärung: das 3'-UTR von BNT162b2 ist nicht erfunden,
sondern aus \textbf{menschlichen} Bausteinen zusammengesetzt -- dem 3'-UTR der
humanen \emph{AES/TLE5}-mRNA und einem Element der mitochondrialen
12S-rRNA (\emph{mtRNR1}). Beides wird unten direkt an der Kassette belegt.

Damit löst sich der Widerspruch auf: \texttt{fs116}--\texttt{fs120} stammen
überhaupt nicht aus dem Spike-Leserahmen. Sie entstehen, wenn das verschobene
Ribosom über das Spike-Stoppcodon hinaus in das 3'-UTR hineinliest -- und dort
übersetzt es \emph{menschliche} Sequenz. Die alte PP2150-Einreichung enthält
denselben \emph{AES}-Abschnitt im selben Raster, deshalb stimmen beide
``Proteine'' überein. Die verbleibenden 12\,\% Unterschied gehen auf die
Abwandlungen des UTR im Impfstoffkonstrukt zurück.

\subsection{Gegenprobe an der Konstruktsequenz}

Diese Erklärung wurde nicht bei der Plausibilität belassen, sondern direkt
geprüft. GenBank führt die Kassette als \texttt{PP544446} (``Vector BNT162b2
clone Vacc\_Batch\_FP9359'', 7692\,bp, hinterlegt 2024); die Spike-CDS liegt
dort auf nt 3578--7399 und übersetzt sich korrekt zu
\texttt{MFVFLVLLPLVS...KLHYT*}, 1273\,aa plus Stopp.

Jedes der 123 Peptide wurde im +1-Rahmen \emph{relativ zum Startcodon} gesucht
und seine Nukleotidposition bestimmt:

\begin{center}
\begin{tabular}{lr}
\toprule
Lage im Konstrukt & Peptide \\
\midrule
innerhalb der Spike-CDS & 113 \\
\textbf{3'-UTR, hinter dem Spike-Stoppcodon} & \textbf{7} \\
5'-UTR, vor dem Startcodon & 1 \\
\midrule
exakt wiedergefunden & 121 von 123 \\
\bottomrule
\end{tabular}
\end{center}

Die sieben Peptide hinter dem Stoppcodon sind \texttt{fs116} bis
\texttt{fs122}, gestaffelt von +5 bis +215\,nt hinter dem Spike-Ende; das eine
davor ist \texttt{fs001} im 5'-UTR. Es sind also \textbf{acht} der 123
Pool-Peptide, die gar nicht aus dem Spike-Leserahmen stammen -- einer mehr als
die fünf, die über den Datenbanktreffer auffielen.

Zuletzt die 3'-Region selbst. Bildet man die 293\,nt hinter dem Stoppcodon auf
die beiden in Frage kommenden Quellen ab, löst sich alles auf:

\begin{center}
\begin{tabular}{llr}
\toprule
3'-Region & Quelle & exakte Übereinstimmung \\
\midrule
nt 10--56 & \emph{TLE5}/\emph{AES} \texttt{NM\_198969.1} ab 1221 & 47\,nt \\
nt 79--145 & \emph{TLE5}/\emph{AES} \texttt{NM\_198969.1} ab 1290 & 67\,nt \\
nt 146--284 & mtRNR1 \texttt{NC\_012920.1} ab 112 & 139\,nt \\
\bottomrule
\end{tabular}
\end{center}

Beide menschlichen Elemente sind damit direkt in der Kassette nachgewiesen und
nicht nur zitiert: vorn das \emph{AES}-Element, dahinter das mitochondriale
12S-Element. Die Lücken zwischen den Blöcken entsprechen den Abwandlungen im
Konstrukt.

Legt man die sieben Peptide auf diese Karte, lässt sich jedes zuordnen:
\texttt{fs116} (+5\,nt) bis \texttt{fs120} (+101\,nt) liegen im
\emph{AES}-Teil -- deshalb treffen sie den \emph{AES}-stämmigen Eintrag
\texttt{AAG23172.1}; \texttt{fs121} (+125\,nt) liegt auf der Nahtstelle, und
\texttt{fs122} (+215\,nt) im mitochondrialen Element. Das achte Peptid,
\texttt{fs001}, liegt vorn im 5'-UTR.

Die Kette ist damit geschlossen und an jedem Glied belegt: Peptide
\texttt{fs116}--\texttt{fs122} liegen hinter dem Spike-Stopp $\rightarrow$ dort
steht das humane \emph{AES/TLE5}-Element $\rightarrow$ dessen +1-Übersetzung
ist zu 88\,\% identisch mit einem menschlichen GenBank-Eintrag, dessen
kodierende Nukleotide ihrerseits zu 150/150 aus \emph{TLE5} stammen.

\subsection*{Zur Genauigkeit der Abschrift}

Die Gegenprobe prüft nebenbei die Abschrift aus dem Tabellenbild mit: 121 der
123 Peptide wurden buchstabengenau in der Kassette wiedergefunden. Bei
\texttt{fs007}, \texttt{fs008}, \texttt{fs039} und \texttt{fs123} fehlt jeweils
ein einzelner Rest zur vollständigen Übereinstimmung (längste vorhandene
Teilsequenz 11 bis 15\,aa). Bei \texttt{fs039} ist das genau die oben
erwähnte, in der Publikation selbst angelegte Inkonsistenz: vorhanden ist
\texttt{PGTGSGSAIAW}, gedruckt ist \texttt{LCTPPGTGSGSAIAW} mit einem
\texttt{P} zu viel. Die übrigen drei dürften Lesefehler aus dem Tabellenbild
sein. Für das Ergebnis ist das ohne Belang: alle vier gehören zur
Spike-Gruppe, die ohnehin keinen Treffer hat.

\subsection{Was daraus für die Ausgangsbehauptung folgt}

Die Formulierung ``völlig neuartige Proteine, meilenweit entfernt vom Spike und
von anderen humanen Proteinen'' trifft zu -- aber nicht überall, und die
Ausnahme ist die interessantere Hälfte:

\begin{itemize}
\item \textbf{Meilenweit vom Spike entfernt: zutreffend, und zwar zwangsläufig.}
      Ein um eine Base verschobener Leserahmen teilt mit dem ursprünglichen
      Protein keinen einzigen Rest; dieselben Nukleotide ergeben in neuer
      Gruppierung völlig andere Codons. Das ist keine Beobachtung, sondern
      Arithmetik.
\item \textbf{Neuartig: zutreffend für 115 der 123 Peptide.} Sie kommen in
      1,18\,Mrd.\ Sequenzen nirgends vor, auch nicht entfernt ähnlich.
\item \textbf{Fern von allem Menschlichen: nicht durchgängig.} Acht der 123
      Peptide stammen gar nicht aus dem Spike-Leserahmen, sieben davon aus dem
      3'-UTR; sie sind im Gegenteil nahezu wörtlich menschlich, weil der
      Impfstoff dort menschliche Bausteine verwendet. Nicht das Ribosom
      erfindet hier etwas, es liest bloß weiter.
\end{itemize}

Zur Vorsicht mahnt noch etwas anderes: diese acht Peptide waren Teil des
ELISpot-Pools. Eine T-Zell-Antwort gegen sie wäre nicht gegen ein
körperfremdes Produkt gerichtet, sondern gegen ein Peptid aus körpereigener
UTR-Sequenz. Das Paper erörtert diesen Punkt nicht; die Einzelantworten je
Peptid wurden nicht erhoben, Pools unterscheiden das nicht.

\subsection{Aminosäurezusammensetzung}

Warum die übrigen Peptide nirgends vorkommen, zeigt ihre Zusammensetzung. Über
alle 123 Peptide (1983 Reste):

\begin{center}
\begin{tabular}{lrlr}
\toprule
Rest & Anteil & Rest & Anteil \\
\midrule
S & 15{,}7\,\% & C & 7{,}3\,\% \\
P & 14{,}8\,\% & R & 6{,}4\,\% \\
T & 13{,}5\,\% & W & 4{,}7\,\% \\
A & 12{,}8\,\% & D+E & 2{,}4\,\% \\
\bottomrule
\end{tabular}
\end{center}

S, P, T, A und C stellen zusammen \textbf{64\,\%} aller Reste; in gewöhnlichen
Proteomen sind es rund 30\,\%. Cystein ist etwa dreifach angereichert,
Tryptophan ebenso, während die sauren Reste Asparaginsäure und Glutaminsäure
auf ein Fünftel des Üblichen fallen und Tyrosin mit 3 von 1983 Resten praktisch
fehlt. Das ist keine Proteinzusammensetzung, sondern niedrigkomplexe,
prolin- und cysteinreiche Sequenz -- die unvermeidliche Folge davon, dass ein
GC-reich codon-optimiertes Gen um eine Base verschoben gelesen wird. Solche
Sequenzen falten schlecht und haben in der Natur kein Gegenstück; dass nr
nichts dazu kennt, ist danach keine Überraschung mehr.

\section{Was dieser Lauf der Studie hinzufügt, und was nicht}

Der Mehrwert gegenüber der Publikation ist klein, und es wäre unredlich, ihn
größer darzustellen. Die Autoren haben den Datenbankabgleich selbst gemacht
(Extended Data Fig.~4). Was hier dazukommt:

\begin{itemize}
\item \textbf{Anderes Werkzeug, größerer Bezugsrahmen.} Das Paper vergleicht
      per BLASTP den BNT162b2-+1-Rahmen gegen den +1-Rahmen des Wildtyp-S-Gens
      -- also zwei Sequenzen gegeneinander. Hier läuft DIAMOND gegen die
      \emph{ganze} nr in ihrem heutigen Stand: 1,18\,Mrd.\ Sequenzen,
      Teil-DBs gebaut am Berichtstag, NCBI-Stand 23.11.2026.
\item \textbf{Nachprüfbares Kriterium statt Abbildung.} Statt eines
      Alignment-Bildes ein maschinell prüfbarer Filter: volle Peptidlänge bei
      100\,\% Identität, ja oder nein.
\item \textbf{Positivkontrolle.} Das Paper führt keine mit. Hier zeigen die
      In-Frame-Luciferase-Peptide, dass der Aufbau Treffer auch findet; ohne
      sie wäre das negative Ergebnis wertlos.
\item \textbf{Eine Frage, die das Paper nicht stellt.} Der Clusterlauf prüft
      nicht nur Identität, sondern entfernte Ähnlichkeit zum menschlichen
      Proteom.
\item \textbf{Ein Befund, der im Paper fehlt.} Acht der 123 Peptide des
      ELISpot-Pools stammen nicht aus dem Spike-Leserahmen, sondern aus den
      UTR des Konstrukts, sieben davon aus dem menschlichen
      \emph{AES/TLE5}-Element im 3'-UTR (Abschnitt~\ref{sec:aehnlich}). Das
      ist mehr als eine Fußnote: eine T-Zell-Antwort gegen diese fünf wäre
      nicht gegen ein fremdes, sondern gegen ein körpereigenes Peptid
      gerichtet. Da die Antworten nur als Pool gemessen wurden, lässt sich
      nachträglich nicht trennen, welcher Anteil der berichteten Reaktivität
      darauf entfällt.
\item \textbf{Laufzeiten und Befehle offengelegt}, siehe Anhang.
\end{itemize}

Was \textbf{nicht} dazukommt: keine neue Biologie, keine Aussage über Risiko
oder Verträglichkeit, keine eigene Messung der Rutschrate, keine klinische
Einordnung. Bestätigt wird ein negatives Ergebnis, das ohnehin zu erwarten war
-- nr enthält keine Mistranslationsprodukte, weil niemand sie einreicht. Der
Beitrag ist Nachprüfbarkeit, nicht Erkenntnis.

\section{Peptidliste: +1FS Spike (Extended Data Table 3)}

\footnotesize
@@peptides@@
\normalsize

\section{Anhang: die ausgeführten Abfragen}

Alles Folgende ist der Stand, der die Zahlen dieses Berichts erzeugt hat.
Skripte und Rohausgaben liegen im Repo \texttt{gerontec/blastDiamond} unter
\texttt{examples/frameshift\_m1psi/}.

\subsection*{A.1 Aufruf}

\begin{lstlisting}
fs_peptide_search.py /tmp/fs_all.fa /tmp/fs_all_out --mode fast
\end{lstlisting}

\subsection*{A.2 Die DIAMOND-Abfrage, je Teil-DB einmal}

\begin{lstlisting}
diamond blastp \
  -d ~/diamond/parts/<part>.dmnd \
  -q /tmp/fs_all.fa \
  -o <out_dir>/<part>.tsv \
  --outfmt 6 qseqid sseqid pident length qlen slen qstart qend \
             sstart send evalue bitscore staxids \
  --evalue 1000 --max-target-seqs 50 \
  --threads 20 --block-size 2 --index-chunks 1 \
  --dbsize @@dbsize@@
\end{lstlisting}

\texttt{-{}-dbsize} ist die Summe der Buchstaben über alle sechs Teile, damit
die E-Werte denen einer Suche gegen die ganze nr entsprechen. Sie stammt aus:

\begin{lstlisting}
SELECT SUM(letters) FROM blast_part;
\end{lstlisting}

Die Empfindlichkeitskontrolle auf \texttt{p1\_mammalia} lief mit denselben
Feldern, aber:

\begin{lstlisting}
  --ultra-sensitive --masking 0 --comp-based-stats 0 --evalue 10000
\end{lstlisting}

\subsection*{A.3 Laufzeitmessung}

Gemessen wird die Wanduhrzeit des Kindprozesses, nicht die Selbstauskunft von
DIAMOND:

\begin{lstlisting}
t0 = time.monotonic()
p  = subprocess.run(cmd, capture_output=True, text=True)
dt = time.monotonic() - t0
\end{lstlisting}

\subsection*{A.4 Exaktheitsfilter}

Aus der Rohausgabe wird nur übernommen, was das ganze Peptid buchstabengenau
abdeckt:

\begin{lstlisting}
if float(c[2]) == 100.0 and int(c[3]) == int(c[4]):   # pident, length == qlen
\end{lstlisting}

\subsection*{A.5 Metadaten-Join gegen die wagodb}

Für jede Accession aus dem Exaktheitsfilter, in Blöcken zu 200:

\begin{lstlisting}
SELECT a.acc, a.taxid, s.len, s.title, t.name, t.rank
  FROM blast_acc  a
  JOIN blast_seq  s ON s.oid   = a.oid
  LEFT JOIN blast_taxon t ON t.taxid = a.taxid
 WHERE a.acc IN (%s, %s, ...);
\end{lstlisting}

Zu beachten: \texttt{blast\_acc} hält eine Zeile je Defline, zusammengefasste
Einträge haben also mehrere Accessions und Taxa je OID. Die von DIAMOND
gemeldete \texttt{staxids} kann daher von der hier verbundenen
\texttt{taxid} abweichen -- genau das ist beim Treffer
\texttt{WP\_212371658.1} der Fall (DIAMOND meldet 10090, die Tabelle führt 2).

\subsection*{A.6 Vorbereitende Abfragen}

Bestand und Zuschnitt der Teil-DBs:

\begin{lstlisting}
SELECT * FROM blast_part;
\end{lstlisting}

Grobe Vorabprüfung, ob überhaupt etwas Einschlägiges in nr steht
(\emph{synthetic construct} 32630, SARS-CoV-2 2697049, Wuhan-Isolat 694009):

\begin{lstlisting}
SELECT taxid, COUNT(*) FROM blast_acc
 WHERE taxid IN (2697049, 32630, 694009) GROUP BY taxid;
-- 32630: 149 433 | 694009: 886 | 2697049: 7 134 543
\end{lstlisting}

Diese Zählung beantwortet die Ausgangsfrage allerdings nicht: Titel und Taxon
sagen nichts darüber, ob eine \emph{Sequenz} vorkommt. Dafür war die
Peptidsuche nötig.

\section{Anhang B: vollständige Reproduktion, Schritt für Schritt}

Alles Folgende ist von Grund auf nachvollziehbar. Benötigt werden nur
\texttt{curl}, \texttt{python3}, \texttt{pdflatex}, DIAMOND, MMseqs2 und ein
Zugang zur nr. Die Skripte liegen im Repo \texttt{gerontec/blastDiamond} unter
\texttt{examples/frameshift\_m1psi/}.

\subsection*{B.1 Quelltexte beschaffen}

\begin{lstlisting}
# Volltext des Papers (JATS-XML)
curl -o pmc.xml "https://eutils.ncbi.nlm.nih.gov/entrez/eutils/efetch.fcgi\
?db=pmc&id=10764286&rettype=xml"

# Extended Data Table 3 (die 123 Peptide) und Table 1 (MS-Peptide) als Bild
B=https://media.springernature.com/full/springer-static/esm/\
art%3A10.1038%2Fs41586-023-06800-3/MediaObjects
curl -o Tab3.jpg "$B/41586_2023_6800_Tab3_ESM.jpg"   # 1398x2521
curl -o Tab1.jpg "$B/41586_2023_6800_Tab1_ESM.jpg"   # 2126x1065
\end{lstlisting}

Die Tabellen liegen ausschließlich als Rasterbild vor; die Peptide mussten
daraus abgeschrieben werden. Die Spalten wurden dafür einzeln vergrößert:

\begin{lstlisting}
convert Tab3.jpg -crop 380x830+618+1690 +repage -resize 300% c3.png
convert Tab3.jpg -crop 420x880+978+0    +repage -resize 300% c4a.png
convert Tab3.jpg -crop 420x880+978+860  +repage -resize 300% c4b.png
\end{lstlisting}

\subsection*{B.2 Abschrift prüfen, bevor sie verwendet wird}

Der Pool ist überlappend aufgebaut, benachbarte Peptide teilen Reste, solange
sie zum selben +1-Abschnitt gehören. Das ergibt eine Selbstkontrolle:

\begin{lstlisting}
check_overlap.py          # 123 Peptide, 8-18 aa, keine Dubletten,
                          # 41 Abbruchstellen = 42 +1-Abschnitte,
                          # 81 Uebergaenge passen buchstabengenau
\end{lstlisting}

\subsection*{B.3 Abfragedatei bauen}

\begin{lstlisting}
{ awk '{printf ">fs%03d_spike1FSpred\n%s\n",      NR, $0}' fs_peptides.txt
  awk '{printf ">ms%02d_fluc1FS_msproven\n%s\n",  NR, $0}' set2.txt
  awk '{printf ">ct%02d_fluc_inframe_control\n%s\n", NR, $0}' set3.txt
} > fs_all.fa          # 123 + 9 + 6 = 138 Sequenzen
\end{lstlisting}

\subsection*{B.4 Suche und Clusterlauf}

\begin{lstlisting}
fs_peptide_search.py  fs_all.fa  out_dir  --mode fast
fs_peptide_search.py  fs_all.fa  out_ultra --mode ultra p1_mammalia   # Kontrolle
fs_peptide_cluster.py fs_all.fa  cluster_out --taxid 9606 --threads 20
best_hits.py          out_dir
\end{lstlisting}

\subsection*{B.5 Den auffälligen Treffer aufklären}

\begin{lstlisting}
# 1. das Protein hinter dem besten Treffer
curl -o aag.gp "https://eutils.ncbi.nlm.nih.gov/entrez/eutils/efetch.fcgi\
?db=protein&id=AAG23172.1&rettype=gp&retmode=text"
#    -> 103 aa, Gen PP2150, "unknown", coded_by AF269289.1:557..868

# 2. die zugehoerige mRNA
curl -o af269289.gb "...efetch.fcgi?db=nuccore&id=AF269289.1&rettype=gb"

# 3. fs116..fs120 auf ihren Ueberlappungen zusammenlegen und ausrichten
align_aag.py        # 50 aa, bester lueckenloser Versatz 36: 44/50 = 88 %

# 4. die kodierenden 150 nt (nt 665..814) gegen TLE5/AES
curl -o tle5.fna "...efetch.fcgi?db=nuccore&id=2462563547,39812018&rettype=fasta"
test_aes.py         # NM_198969.1: 150 von 150 nt identisch ab Position 1216
\end{lstlisting}

Achtung bei Schritt 4: \texttt{NM\_005938} ist \emph{nicht} AES, sondern
FOXO4. Die richtige Kennung findet man über
\texttt{esearch.fcgi?db=nuccore\&term=TLE5[gene]+AND+"Homo sapiens"[orgn]\
+AND+biomol\_mrna[prop]+AND+srcdb\_refseq[prop]}.

\subsection*{B.6 Gegenprobe an der Konstruktsequenz}

\begin{lstlisting}
curl -o bnt.gb "...efetch.fcgi?db=nuccore&id=PP544446&rettype=gb&retmode=text"
verify_utr.py       # Spike-CDS 3578..7399, +1-Rahmen ab Startcodon
                    # 113 Peptide in der CDS, 7 im 3'-UTR, 1 im 5'-UTR
\end{lstlisting}

Zwei Fallstricke, die beim Nachbauen Zeit kosten:

\begin{enumerate}
\item Beim Einlesen der GenBank-Datei dürfen aus dem \texttt{ORIGIN}-Block
      \emph{alle} Buchstaben übernommen werden, nicht nur \texttt{ACGTN}.
      Filtert man IUPAC-Mehrdeutigkeitscodes weg, fehlen hier fünf Basen
      (7687 statt 7692\,bp) und der Leserahmen ist verschoben -- die CDS
      übersetzt sich dann zu \texttt{SCSWCCCLWCPA...} statt
      \texttt{MFVFLVLLPLVS...}, was als Fehler leicht zu übersehen ist, weil
      es wie ein +1-Produkt aussieht.
\item Der gesuchte +1-Rahmen ist der gegenüber dem \emph{Startcodon}
      verschobene, nicht der gegenüber dem Vektoranfang. Mit dem falschen
      Bezugspunkt findet man null von 123 Peptiden.
\end{enumerate}

\subsection*{B.7 Zusammensetzung und Bericht}

\begin{lstlisting}
composition.py                       # S+P+T+A+C = 64 % der Reste
fs_report_pdf.py out_dir fs_all.fa Frameshift.pdf cluster_out
\end{lstlisting}

\subsection*{B.8 Verwendete Kennungen auf einen Blick}

\begin{center}
\footnotesize
\begin{tabular}{lll}
\toprule
Kennung & Was & Wozu \\
\midrule
PMC10764286 / PMID 38057663 & Mulroney et al., Nature 625 & Quelle \\
PMID 38867495 & Boros et al., Review & angebliche Folgestudie \\
\texttt{PP544446} & BNT162b2-Kassette, 7692\,bp & Gegenprobe \\
\texttt{AAG23172.1} / \texttt{AF269289.1} & PP2150, ``unknown'' & bester Treffer \\
\texttt{NM\_198969.1} & TLE5 / AES mRNA & Herkunft des 3'-UTR \\
\texttt{NC\_012920.1}:648--1601 & mtRNR1 / 12S & zweites UTR-Element \\
\texttt{MSV000093074}, \texttt{GSE223044} & MS- und RNA-seq-Daten & Hinterlegung \\
\texttt{P08659} & Firefly-Luciferase & Positivkontrolle \\
\bottomrule
\end{tabular}
\end{center}

\end{document}
